\chapter{The Trouble with Defining Intelligence by Discovery}
\label{ch:discovery-trouble}

\begin{plainlanguage}
\textbf{In plain terms:} Every generation confidently redefines intelligence, and every next generation finds the definition wrong. That cycling, rather than converging, is itself evidence that intelligence isn't a single hidden thing waiting to be correctly named. This chapter states that puzzle precisely; it doesn't yet solve it.
\end{plainlanguage}

Ask a roomful of researchers to define intelligence and you will get a roomful of different answers, and this would not be troubling if the answers were converging --- if each decade's definition were a refinement of the last, narrowing in on some fixed target the way successive measurements of a physical constant narrow in on its true value. They are not converging. They are cycling.

Intelligence has been defined, at various points and by serious people, as the capacity to optimize a reward signal, as the capacity to predict future sensory states, as the capacity to compress experience into efficient representations, as a form of computation, as the manipulation of symbols according to rules, as an emergent property of sufficiently complex information processing, as adaptive behavior in a changing environment, and as whatever a particular test happens to measure. Each definition arrived with the confidence of a discovery --- \emph{this} is what intelligence actually is, beneath the folk confusion --- and each was eventually absorbed, qualified, or quietly abandoned by the next generation, which then proposed its own discovery with the same confidence.

This pattern deserves a name, because once named it becomes hard to unsee: call it the \emph{discovery cycle}. A theory is proposed as if intelligence were a fixed target being progressively triangulated. The theory illuminates real cases --- usually the cases its proponents had in mind when constructing it --- and fails on cases nobody had in mind. The failures accumulate. A new theory is proposed, again as a discovery, again illuminating a different set of cases and failing on others. Nothing about this process looks like convergence on a stable target. It looks like something else.

The most natural explanation --- and the one this book will spend its length defending, complicating, and ultimately earning the right to state precisely --- is that ``intelligence'' was never a fixed target in the first place. Not because there is nothing real underneath the word, but because what is underneath the word is not the kind of thing a single definition is built to capture: not a latent variable waiting to be measured more precisely, but a structured, multidimensional pattern that different observers, working from different vantage points and against different failure cases, partially and unevenly perceive. Each definition is not wrong so much as a true description of one face of something that has several.

If that is right, then the discovery cycle is not a series of failures to find the right definition. It is the visible signature of a different process entirely --- one in which a concept gets \emph{built}, not found, through repeated contact between candidate descriptions and the world's resistance to them. This book calls that process \emph{perspectival convergence}, and the rest of Part~I exists to state the problem precisely enough that the machinery built in Part~II has something exact to do.

\section*{Notation}

Let
\[
\Theta = \{\theta_1, \theta_2, \ldots\}
\]
denote the space of candidate definitions of intelligence --- every precise proposal that has been or could be advanced, from ``intelligence is the capacity to maximize expected reward'' to ``intelligence is whatever an IQ test measures'' to far more recent and far stranger candidates. Equip $\Theta$ with a conceptual-distance metric $d_\Theta(\theta_i,\theta_j)$, a rough and for now informal measure of how far apart two candidate definitions are in the space of things they would classify as intelligent or not.

The discovery cycle, restated in this notation, is the observation that the historically realized sequence of candidate definitions does not look like a sequence converging toward a single $\theta^*$. It looks, across the twentieth and early twenty-first centuries, like $|\Theta_t|$ --- the number of seriously defended candidates in active circulation --- has tended to grow rather than shrink, even as individual theories have died.

\begin{conjecture}[Definition-Space Expansion]
In the historical record of intelligence research, $|\Theta_t|$ increases for long stretches before any convergence mechanism causes it to decrease.
\end{conjecture}

This is stated as a conjecture rather than an established fact because it has not, in this book, been checked against an actual census of the literature --- it is offered as the intuition that motivates everything that follows, and a reader who suspects it is false (perhaps definitions really have been narrowing, just slowly, and the appearance of expansion is an artifact of which decades get remembered) should treat the rest of the book as conditional on it being roughly right. What the conjecture is for is this: if true, it means the field needs an elimination mechanism --- some way of taking the expanding set $\Theta_t$ and producing a narrower, more defensible $\Theta_T$ --- and it means that mechanism cannot simply be ``wait for more research,'' since more research, on the historical evidence, has been expanding $\Theta$ rather than shrinking it. Chapter~\ref{ch:candidate-space} builds the first candidate for such a mechanism. It will turn out, in Chapter~\ref{ch:nonmonotonic}, to be the wrong one.
